% Source: adaboost.tex
\begin{notechapter}
\chapter{Ensemble Methods}\label{ch:ensemble-methods}
\chaptersubtitle{AdaBoost, forward stagewise additive modeling and gradient boosting}

\section{AdaBoost Algorithm}

AdaBoost (Adaptive Boosting) is an ensemble learning algorithm that 
combines multiple weak learners (e.g., shallow decision trees) to 
create a strong learner. It focuses on instances that are difficult 
to classify, iteratively improving the model's performance.

\subsection{How AdaBoost Works}

\subsubsection{Initialization}
\begin{itemize}
    \item Start with a dataset of \( N \) training instances: 
    \( (x_1, y_1), (x_2, y_2), \dots, (x_N, y_N) \), where \( x_i \) 
    are features and \( y_i \) are labels (e.g., \( y_i \in \{-1, 1\} \)).
    \item Assign equal weights to all instances: 
    \( w_i = \frac{1}{N} \) for \( i = 1, 2, \dots, N \).
\end{itemize}

\subsubsection{Iterative Training}
AdaBoost trains \( T \) weak learners. For each iteration 
\( t = 1, 2, \dots, T \):

\begin{itemize}
    \item \textbf{Step 1: Train a Weak Learner}
    \begin{itemize}
        \item Train a weak learner \( h_t(x) \) on the weighted data.
        \item Compute the weighted error:
        \[
        \epsilon_t = \sum_{i=1}^N w_i \cdot \mathbb{I}(h_t(x_i) \neq y_i),
        \]
        where \( \mathbb{I} \) is 1 if incorrect, 0 otherwise.
    \end{itemize}

    \item \textbf{Step 2: Compute Weak Learner's Weight}
    \begin{itemize}
        \item Compute \( \alpha_t \) for \( h_t(x) \):
        \[
        \alpha_t = \frac{1}{2} \ln\left(\frac{1 - \epsilon_t}{\epsilon_t}\right).
        \]
        \item If \( \epsilon_t \) is small, \( \alpha_t \) is large.
        \item If \( \epsilon_t \) is large, \( \alpha_t \) is small.
    \end{itemize}

    \item \textbf{Step 3: Update Instance Weights}
    \begin{itemize}
                \item Update the weights, with labels and predictions in \( \{-1, +1\} \):
        \[
        w_i \leftarrow w_i \cdot \exp\left(-\alpha_t \, y_i \, h_t(x_i)\right),
        \]
        i.e.\ multiply by \( e^{\alpha_t} > 1 \) the weights of misclassified instances
        (\( y_i h_t(x_i) = -1 \)) and by \( e^{-\alpha_t} < 1 \) those of correctly classified ones.
        Up to the normalization below, this is the same as multiplying only the misclassified weights by
        \( e^{2\alpha_t} = \frac{1-\epsilon_t}{\epsilon_t} \). After normalization, the misclassified instances
        carry exactly half of the total weight, so \( h_t \) has weighted error \( 1/2 \) under the new weights.
        \item Normalize weights so they sum to 1:
        \[
        w_i \leftarrow \frac{w_i}{\sum_{j=1}^N w_j}.
        \]
    \end{itemize}
\end{itemize}

\subsubsection{Combine Weak Learners}
After \( T \) iterations, the final strong learner \( H(x) \) is:
\[
H(x) = \text{sign}\left(\sum_{t=1}^T \alpha_t \cdot h_t(x)\right).
\]
\begin{itemize}
    \item Weak learners with lower error contribute more.
    \item The \( \text{sign} \) function outputs binary predictions.
\end{itemize}

\subsection{Key Features of AdaBoost}
\begin{itemize}
    \item Focuses on hard-to-classify instances.
    \item Adaptively weights weak learners based on performance.
    \item Reduces bias and improves generalization.
    \item Less prone to overfitting with simple weak learners.
\end{itemize}

\subsection{Advantages of AdaBoost}
\begin{itemize}
    \item Simple and easy to implement.
    \item Works well with weak learners like decision stumps.
    \item Achieves high accuracy with minimal tuning.
\end{itemize}

\subsection{Limitations of AdaBoost}
\begin{itemize}
    \item Sensitive to noisy data and outliers.
    \item Requires weak learners to perform better than random guessing.
\end{itemize}

\section{Forward Stagewise Additive Modeling (FSAM)}

FSAM is a framework for building additive models sequentially. It is 
the theoretical foundation for many boosting algorithms, including 
AdaBoost and Gradient Boosting.

\subsection{How FSAM Works}
1. \textbf{Additive Model}:
   The goal is to build a model \( F(x) \) as a sum of simpler models:
   \[
   F(x) = \sum_{t=1}^T \alpha_t h_t(x),
   \]
   where \( h_t(x) \) is a weak learner and \( \alpha_t \) is its weight.

2. \textbf{Sequential Training}:
   \begin{itemize}
       \item At each step \( t \), a new weak learner \( h_t(x) \) is added.
       \item The weak learner is chosen to minimize a loss function \( L \):
         \[
         h_t(x) = \arg\min_{h} \sum_{i=1}^N L\left(y_i, F_{t-1}(x_i) + \alpha_t h(x_i)\right).
         \]
       \item The weight \( \alpha_t \) is determined via optimization.
   \end{itemize}

3. \textbf{Greedy Optimization}:
   \begin{itemize}
       \item FSAM is greedy because it optimizes one weak learner at a time.
   \end{itemize}

\subsection{Connection to AdaBoost}
AdaBoost is a special case of FSAM where:
\begin{itemize}
    \item The loss function is exponential: 
    \[
    L(y, F(x)) = \exp(-y F(x)).
    \]
    \item The weak learners are binary classifiers (e.g., decision stumps).
\end{itemize}

\section{Gradient Boosting}

Gradient Boosting is a generalization of boosting algorithms that 
uses gradient descent to minimize a differentiable loss function.

\subsection{How Gradient Boosting Works}
1. \textbf{Additive Model}:
   The model is built sequentially:
   \[
   F(x) = \sum_{t=1}^T h_t(x),
   \]
   where \( h_t(x) \) is a weak learner (often a decision tree).

2. \textbf{Gradient Descent in Function Space}:
   \begin{itemize}
       \item At each step \( t \), compute pseudo-residuals:
         \[
         r_{i,t} = -\frac{\partial L(y_i, F_{t-1}(x_i))}{\partial F_{t-1}(x_i)}.
         \]
       \item Train a weak learner \( h_t(x) \) to predict \( r_{i,t} \):
         \[
         h_t(x) = \arg\min_{h} \sum_{i=1}^N (r_{i,t} - h(x_i))^2.
         \]
   \end{itemize}

3. \textbf{Update the Model}:
   \begin{itemize}
       \item Update the model:
         \[
         F_t(x) = F_{t-1}(x) + \nu h_t(x),
         \]
         where \( \nu \) is the learning rate.
   \end{itemize}

4. \textbf{Iterate Until Convergence}:
   \begin{itemize}
       \item Repeat for \( T \) iterations or until the loss stops improving.
   \end{itemize}

\subsection{Connection to FSAM}
Gradient Boosting is a specific implementation of FSAM where:
\begin{itemize}
    \item The loss function is any differentiable loss.
    \item The weak learners are trained to approximate the negative gradient.
\end{itemize}

\subsection{Example: Gradient Boosting for Regression}
For regression with squared error loss:
\begin{itemize}
    \item Pseudo-residuals:
      \[
      r_{i,t} = y_i - F_{t-1}(x_i).
      \]
    \item Weak learner: A decision tree predicts \( r_{i,t} \).
    \item Update model:
      \[
      F_t(x) = F_{t-1}(x) + \nu h_t(x).
      \]
\end{itemize}

\section{Key Differences Between FSAM and Gradient Boosting}

\begin{tabular}{|l|l|l|}
\hline
\textbf{Feature} & \textbf{FSAM} & \textbf{Gradient Boosting} \\
\hline
Loss Function & Any loss function & Differentiable loss function \\
\hline
Weak Learner & Any weak learner & Typically decision trees \\
\hline
Optimization & Greedy, stepwise & Gradient descent in function space \\
\hline
Weight Calculation & Analytically derived & Determined via line search \\
\hline
Flexibility & Less flexible & Highly flexible \\
\hline
\end{tabular}

\section{Advantages of Gradient Boosting}
\begin{itemize}
    \item Flexibility: Can handle any differentiable loss function.
    \item Powerful: Often achieves state-of-the-art performance.
    \item Interpretability: Weak learners (e.g., trees) are interpretable.
\end{itemize}

\section{Limitations of Gradient Boosting}
\begin{itemize}
    \item Computationally expensive: Sequential training can be slow.
    \item Sensitive to overfitting: Requires careful tuning.
    \item Black-box nature: The ensemble can be complex.
\end{itemize}
\end{notechapter}
